% ------------------------------------------------------------------------
% file `3_uaa2_2022-connaitre-1-exercise-body.tex'
%   in folder `xsim/'
%
%     exercise of type `connaitre' with id `1'
%
% generated by the `connaitre' environment of the
%   `xsim' package v0.21 (2022/02/12)
% from source `3_uaa2_2022' on 2023/03/15 on line 25
% ------------------------------------------------------------------------
    Vrai ou faux ? Justifie dans les deux cas.
    \begin{enumerate}
        \item Deux nombres opposés ont la même racine carrée.
        \item \(\sqrt{20}=10\sqrt{2}\).
        \item Si \(ABC\) est un triangle rectangle en \(C\) alors \(\cos\left(\widehat{A}\right) = \dfrac{|AC|}{|BC|}\).
        \item Si \(ABC\) est un triangle rectangle en \(C\) alors \(\cos^2\left(\widehat{A}\right) + \sin^2\left(\widehat{A}\right)=1\).
        \item Dans tout triangle rectangle, la tangente des angles aigu est le rapport entre leur cosinus et leur sinus.
        \item Dans un triangle rectangle isocèle, le sinus et le cosinus d'un même angle sont égaux.
        \item \(\sqrt{-16}=4\).
        \item Dans un triangle rectangle de côtés 10, 3 et \(\sqrt{91}\), un des angles mesure \ang{17,46} (l'angle est arrondi au centième).
    \end{enumerate}
